\documentclass[12pt,letterpaper]{article}
\usepackage{graphicx,textcomp}
\usepackage{natbib}
\usepackage{setspace}
\usepackage{fullpage}
\usepackage{color}
\usepackage[reqno]{amsmath}
\usepackage{amsthm}
\usepackage{fancyvrb}
\usepackage{amssymb,enumerate}
\usepackage[all]{xy}
\usepackage{endnotes}
\usepackage{lscape}
\newtheorem{com}{Comment}
\usepackage{float}
\usepackage{hyperref}
\newtheorem{lem} {Lemma}
\newtheorem{prop}{Proposition}
\newtheorem{thm}{Theorem}
\newtheorem{defn}{Definition}
\newtheorem{cor}{Corollary}
\newtheorem{obs}{Observation}
\usepackage[compact]{titlesec}
\usepackage{dcolumn}
\usepackage{tikz}
\usetikzlibrary{arrows}
\usepackage{multirow}
\usepackage{xcolor}
\newcolumntype{.}{D{.}{.}{-1}}
\newcolumntype{d}[1]{D{.}{.}{#1}}
\definecolor{light-gray}{gray}{0.65}
\usepackage{url}
\usepackage{listings}
\usepackage{color}

\definecolor{codegreen}{rgb}{0,0.6,0}
\definecolor{codegray}{rgb}{0.5,0.5,0.5}
\definecolor{codepurple}{rgb}{0.58,0,0.82}
\definecolor{backcolour}{rgb}{0.95,0.95,0.92}

\lstdefinestyle{mystyle}{
	backgroundcolor=\color{backcolour},   
	commentstyle=\color{codegreen},
	keywordstyle=\color{magenta},
	numberstyle=\tiny\color{codegray},
	stringstyle=\color{codepurple},
	basicstyle=\footnotesize,
	breakatwhitespace=false,         
	breaklines=true,                 
	captionpos=b,                    
	keepspaces=true,                 
	numbers=left,                    
	numbersep=5pt,                  
	showspaces=false,                
	showstringspaces=false,
	showtabs=false,                  
	tabsize=2
}
\lstset{style=mystyle}
\newcommand{\Sref}[1]{Section~\ref{#1}}
\newtheorem{hyp}{Hypothesis}

\title{Problem Set 2}
\date{Due: October 23, 2025}
\author{Applied Stats/Quant Methods 1}

\begin{document}
	\maketitle
	\section*{Instructions}
\begin{itemize}
	\item Please show your work! You may lose points by simply writing in the answer. If the problem requires you to execute commands in \texttt{R}, please include the code you used to get your answers. Please also include the \texttt{.R} file that contains your code. If you are not sure if work needs to be shown for a particular problem, please ask.
	\item Your homework should be submitted electronically on GitHub.
	\item This problem set is due before 23:59 on Thursday October 23, 2025. No late assignments will be accepted.

\end{itemize}

	
	\vspace{.5cm}
	\section*{Question 1: Political Science}
		\vspace{.25cm}
	The following table was created using the data from a study run in a major Latin American city.\footnote{Fried, Lagunes, and Venkataramani (2010). ``Corruption and Inequality at the Crossroad: A Multimethod Study of Bribery and Discrimination in Latin America. \textit{Latin American Research Review}. 45 (1): 76-97.} As part of the experimental treatment in the study, one employee of the research team was chosen to make illegal left turns across traffic to draw the attention of the police officers on shift. Two employee drivers were upper class, two were lower class drivers, and the identity of the driver was randomly assigned per encounter. The researchers were interested in whether officers were more or less likely to solicit a bribe from drivers depending on their class (officers use phrases like, ``We can solve this the easy way'' to draw a bribe). The table below shows the resulting data.

\newpage
\begin{table}[h!]
	\centering
	\begin{tabular}{l | c c c }
		& Not Stopped & Bribe requested & Stopped/given warning \\
		\\[-1.8ex] 
		\hline \\[-1.8ex]
		Upper class & 14 & 6 & 7 \\
		Lower class & 7 & 7 & 1 \\
		\hline
	\end{tabular}
\end{table}

\begin{enumerate}
	
	\item [(a)]
	Calculate the $\chi^2$ test statistic by hand/manually (even better if you can do "by hand" in \texttt{R}).\\
	
	 \textbf{Answer:} The chi-square test statistic is:
	\[
	\chi^2 = 3.791, \quad df = 2
	\]
	Computed in \texttt{R} as:
	\lstinputlisting[language=R, firstline=37, lastline=67]{PS02_JL.R}
	
	\item [(b)]
	Now calculate the p-value from the test statistic you just created (in \texttt{R}).\footnote{Remember frequency should be $>$ 5 for all cells, but let's calculate the p-value here anyway.}  What do you conclude if $\alpha = 0.1$?\\
	
	\textbf{Answer:} The p-value obtained was $p = 0.150$, which is greater than $\alpha = 0.1$. 
	We therefore \textit{fail to reject} the null hypothesis and conclude that there is no statistically significant association between social class and outcome of the stop.\\
	
	Computed in \texttt{R} as:
	\lstinputlisting[language=R, firstline=71, lastline=78]{PS02_JL.R}
	
	\newpage
	\item [(c)] Calculate the standardized residuals for each cell and put them in the table below.
	\vspace{1cm}


\begin{table}[h]
	\centering
	\begin{tabular}{l | c c c }
		& Not Stopped & Bribe requested & Stopped/given warning \\
		\\[-1.8ex] 
		\hline \\[-1.8ex]
		Upper class  & 0.136 & -0.815 & 0.819 \\
		\\
		Lower class  & -0.183 & 1.094 & -1.099 \\
	\end{tabular}
\end{table}
	\vspace{1em}
	Computed in \texttt{R} as:
	\lstinputlisting[language=R, firstline=81, lastline=83]{PS02_JL.R}
	\vspace{2cm}
	
	\item [(d)] How might the standardized residuals help you interpret the results?  
	
	\textbf{Answer:} 
	If the Chi-square test were significant (which it was not in this case), we would know that the variables are related, but not which specific categories drive that relationship. The standardised residuals address this by showing how much each observed count deviates from its expected count, measured in units of standard deviation. 
	
	Values close to zero indicate that the observed count is very similar to what would be expected under the null hypothesis of independence, whereas residuals greater than $\lvert 2 \rvert$ are typically considered large enough to indicate a significant deviation. In our case, none of the standardised residuals exceed $\lvert 2 \rvert$, which confirms the earlier finding that the Chi-square test was not statistically significant at $\alpha = 0.10$.
	
\end{enumerate}
\newpage

\section*{Question 2: Economics}
Chattopadhyay and Duflo were interested in whether women promote different policies than men.\footnote{Chattopadhyay and Duflo. (2004). ``Women as Policy Makers: Evidence from a Randomized Policy Experiment in India. \textit{Econometrica}. 72 (5), 1409-1443.} Answering this question with observational data is pretty difficult due to potential confounding problems (e.g. the districts that choose female politicians are likely to systematically differ in other aspects too). Hence, they exploit a randomized policy experiment in India, where since the mid-1990s, $\frac{1}{3}$ of village council heads have been randomly reserved for women. A subset of the data from West Bengal can be found at the following link: \url{https://raw.githubusercontent.com/kosukeimai/qss/master/PREDICTION/women.csv}\\

\noindent Each observation in the data set represents a village and there are two villages associated with one GP (i.e. a level of government is called "GP"). Figure~\ref{fig:women_desc} below shows the names and descriptions of the variables in the dataset. The authors hypothesize that female politicians are more likely to support policies female voters want. Researchers found that more women complain about the quality of drinking water than men. You need to estimate the effect of the reservation policy on the number of new or repaired drinking water facilities in the villages.
\vspace{.5cm}
%\begin{figure}[h!] %graphic missing%
	%\caption{\footnotesize{Names and description of variables from Chattopadhyay and Duflo (2004).}}
	%\vspace{.5cm}
	%\centering
	%\label{fig:women_desc}
	%\includegraphics[width=1.1\textwidth]{../../../../graphics/women_desc.png}  
%\end{figure}		

\newpage
\begin{enumerate}
	\item [(a)] State a null and alternative (two-tailed) hypothesis. 
	
	\textbf{Answer:} 
	The null hypothesis assumes that there is \textit{no difference} between villages with reserved seats for women and those without. 
	The alternative hypothesis assumes that the reservation policy \textit{does have an effect}, meaning there is a difference, 
	but it does not specify the direction of that difference.

	\[
	H_0: \ \mu_{\text{reserved}} = \mu_{\text{not reserved}}
	\quad \text{vs.} \quad
	H_a: \ \mu_{\text{reserved}} \neq \mu_{\text{not reserved}}
	\]
		
	\item [(b)] Run a bivariate regression to test this hypothesis in \texttt{R} (include your code!).
	
	\textbf{Answer:} 
	A bivariate regression was estimated to test whether villages with reserved seats for women differ from those without in the outcome variable.
	
	% Table created by stargazer v.5.2.3 by Marek Hlavac, Social Policy Institute. E-mail: marek.hlavac at gmail.com
	% Date and time: Tue, Oct 14, 2025 - 09:45:08
	\begin{table}[!htbp] \centering 
		\begin{tabular}{@{\extracolsep{5pt}}lc} 
			\\[-1.8ex]\hline 
			\hline \\[-1.8ex] 
			& \multicolumn{1}{c}{\textit{Dependent variable:}} \\ 
			\cline{2-2} 
			\\[-1.8ex] & Outcome \\ 
			\hline \\[-1.8ex] 
			Reserved (1 = Yes) & 9.252$^{**}$ \\ 
			& (3.948) \\ 
			Constant & 14.738$^{***}$ \\ 
			& (2.286) \\ 
			\hline \\[-1.8ex] 
			Observations & 322 \\ 
			R$^{2}$ & 0.017 \\ 
			Adjusted R$^{2}$ & 0.014 \\ 
			Residual Std. Error & 33.446 (df = 320) \\ 
			F Statistic & 5.493$^{**}$ (df = 1; 320) \\ 
			\hline 
			\hline \\[-1.8ex] 
			\textit{Note:}  & \multicolumn{1}{r}{$^{*}$p$<$0.1; $^{**}$p$<$0.05; $^{***}$p$<$0.01} \\ 
		\end{tabular} 
	\end{table} 
	\vspace{1em}
	Computed in \texttt{R} as:
	\lstinputlisting[language=R, firstline=100, lastline=122]{PS02_JL.R}
	
	\vspace{1cm}
	\item [(c)] Interpret the coefficient estimate for reservation policy. 
	
	\textbf{Answer:} 
	The estimated intercept, $\hat{\beta}_0 = 14.74$, represents the average number of new or repaired water facilities in villages where the seat was \textit{not} reserved for women. The coefficient for the reservation policy, $\hat{\beta}_1 = 9.252$, captures the estimated causal effect of the policy. It indicates that villages with a reserved seat had, on average, about 9.25 more new or repaired water facilities than those without a reserved seat. 
	
	The p-value associated with $\hat{\beta}_1$ is $p = 0.0197$, which is below $\alpha = 0.05$. We therefore \textit{reject the null hypothesis} that the reservation policy had no effect. The policy is associated with a statistically significant increase in the number of water facilities. 
	
	However, the $R^2$ value of $0.0169$ suggests that the reservation policy explains only about 1.7\% of the total variance in the outcome. While statistically significant, the effect size is small relative to the overall variation in the data.
	
\end{enumerate}

%\newpage
%	\section*{Question 3 (30 points): Biology}
%
%There is a physiological cost of reproduction for fruit flies, such that it reduces the lifespan of female fruit flies.  Is there a similar cost to male fruit flies?  This dataset contains observations from five groups of 25 male fruit flies. The experiment tests if increased reproduction reduces longevity for male fruit flies. The five groups are: males forced to live alone, males assigned to live with one or eight newly pregnant females (non-receptive females), and males assigned to live with one or eight virgin females (interested females). The name of the data set is \texttt{fruitfly.csv}.\footnote{Partridge and Farquhar (1981).``Sexual Activity and the Lifespan of Male Fruitflies''. \textit{Nature}. 294, 580-581.}
%	\vspace{1cm}
%
%\begin{tabular}{r|l}
%	\texttt{No} & serial number (1-25) within each group of 25\\
%	\texttt{type} & Type of experimental assignment \\
%	& \hspace{0.1in} $1=$ no females  \\
%	& \hspace{0.1in} $2=$ 1 newly pregnant female \\
%	& \hspace{0.1in} $3=$ 8 newly pregnant females\\
%	& \hspace{0.1in} $4=$ 1 virgin female\\
%	& \hspace{0.1in} $5=$ 8 virgin females\\
%	\texttt{lifespan} & lifespan (days)\\
%	\texttt{thorax} & length of thorax (mm)\\
%	\texttt{sleep} & percentage of each day spent sleeping\\
%\end{tabular}
%	\vspace{1cm}
%\begin{enumerate}
%	
%	\item
%	Import the data set and obtain summary statistiscs and examine the distribution of the overall lifespan of the fruitflies.  
%
%\newpage
%	\item
%	Plot \texttt{lifespan} vs \texttt{thorax}. Does it look like there is a linear relationship? Provide the plot. What is the correlation coefficient between these two variables?
%		\vspace{6cm}
%	\item
%	Regress \texttt{lifespan} on \texttt{thorax}.  Interpret the slope of the fitted model.
%			\vspace{6cm}
%	\item
%	Test for a significant linear relationship between  \texttt{lifespan} and \texttt{thorax}. Provide and interpret your results of your test.
%	
%\newpage
%	\item
%	
%	Provide the 90\% confidence interval for the slope of the fitted model.
%	
%			\vspace{.5cm}
%	\begin{itemize}
%		\item
%		Use the formula of confidence interval.		\vspace{.5cm}
%		\item
%		Use the function  \texttt{confint()}  in \texttt{R} .
%	\end{itemize}
%			\vspace{6cm}
%	\item Use the \texttt{predict()} function in \texttt{R} to (1) predict an individual fruitfly's lifespan when \texttt{thorax}=0.8 and (2) the average \texttt{lifespan} of fruitflies when \texttt{thorax}=0.8 by the fitted model. This requires that you compute prediction and confidence intervals. What are the expected values of lifespan? What are the prediction and confidence intervals around the expected values? 
%	
%			\vspace{6cm}
%	\item	For a sequence of \texttt{thorax} values, draw a plot with their fitted values for \texttt{lifespan}, as well as the prediction intervals and confidence intervals.
%


%\end{enumerate}
\end{document}
